% GENERATED FILE. Edit canonical lessons, then run npm run generate:books.
% canonical-source-hash: fnv1a64:839004b0f0208248
% canonical-lessons: BN-C317-bir, BN-C317-atmotyag, BN-C317-shohid, BN-C317-atmosomorpon-kora, BN-C317-porajoy

\chapter{Courage and Sacrifice}
\label{ch:bn-courage-and-sacrifice}
\hlchaptermodality{317}

\begin{chapteropening}
\textbf{By the end of this chapter:} \emph{I can talk about heroes, martyrs, surrender and defeat.}

Complete the last lesson of chapter 317: I can talk about heroes, martyrs, surrender and defeat.
\end{chapteropening}

\section[\texorpdfstring{\bn{বীর} (bīr)}{বীর (bīr)}]{\bn{বীর} (bīr) — a hero; a warrior}

\label{lesson:BN-C317-bir}



\begin{quote}
Before the new one: say the Bengali for to hold out, then the Bengali for to fall back.
\end{quote}

\subsection*{You'll want to know: \bn{বীর}}
\textbf{\bn{বীর}} — \emph{bīr} — \textquotedblleft{}a hero; a warrior\textquotedblright{}.

\begin{quote}
\bn{সেই} \bn{যুদ্ধের} \bn{বীরদের} \bn{গল্প} \bn{সবাই} \bn{জানে।}
\end{quote}

\begin{quote}
\emph{sei juddher bīrder gôlpo sôbāi jāne.}
\end{quote}

\begin{quote}
Everyone knows the stories of the heroes of that war.
\end{quote}

\begin{cousinweb}
\textbf{\bn{বীর}} (\emph{bīr}) is Sanskrit \emph{vīra}, 'a man; a hero, a chief' (Monier-Williams); the \bn{ী} is long in writing.
\end{cousinweb}

\begin{grammarlens}[title={What you've built}]
One more word for telling what happened.
\end{grammarlens}

\subsection*{Your turn}
\begin{itemize}
\raggedright
  \item \emph{Say it:} \emph{bīr}
  \item \emph{Say it:} \emph{bīr}, once more
  \item \emph{Say it:} \emph{sei juddher bīrder gôlpo sôbāi jāne.} — Everyone knows the stories of the heroes of that war.
  \item \emph{Recall:} say the Bengali for to hold out, then the Bengali for to fall back, then say \emph{bīr} again
\end{itemize}

\begin{tcolorbox}[breakable,colback=teal!4,colframe=teal!35!black,title={Before you move on}]
What does \bn{বীর} mean? (\textquotedblleft{}A hero; a warrior\textquotedblright{}.) Say it once more.
\end{tcolorbox}
\section[\texorpdfstring{\bn{আত্মত্যাগ} (ātmôtyāg)}{আত্মত্যাগ (ātmôtyāg)}]{\bn{আত্মত্যাগ} (ātmôtyāg) — self-sacrifice}

\label{lesson:BN-C317-atmotyag}



\begin{quote}
Before the new one: say the Bengali for to fall back, then the Bengali for a warrior.
\end{quote}

\subsection*{You'll want to know: \bn{আত্মত্যাগ}}
\textbf{\bn{আত্মত্যাগ}} — \emph{ātmôtyāg} — \textquotedblleft{}self-sacrifice\textquotedblright{}.

\begin{quote}
\bn{দেশের} \bn{জন্য} \bn{সৈন্যদের} \bn{আত্মত্যাগ} \bn{আমরা} \bn{মনে} \bn{রাখব।}
\end{quote}

\begin{quote}
\emph{desher jônyô soinyôder ātmôtyāg āmrā mône rākhbô.}
\end{quote}

\begin{quote}
We will remember the soldiers' sacrifice for their country.
\end{quote}

\begin{cousinweb}
\textbf{\bn{আত্মত্যাগ}} (\emph{ātmôtyāg}) joins Sanskrit \emph{ātman}, 'self', to \emph{tyāga}, 'giving up, abandoning' (Monier-Williams): giving up oneself.
\end{cousinweb}

\begin{grammarlens}[title={What you've built}]
One more word for telling what happened.
\end{grammarlens}

\subsection*{Your turn}
\begin{itemize}
\raggedright
  \item \emph{Say it:} \emph{ātmôtyāg}
  \item \emph{Say it:} \emph{ātmôtyāg}, once more
  \item \emph{Say it:} \emph{desher jônyô soinyôder ātmôtyāg āmrā mône rākhbô.} — We will remember the soldiers' sacrifice for their country.
  \item \emph{Recall:} say the Bengali for to fall back, then the Bengali for a warrior, then say \emph{ātmôtyāg} again
\end{itemize}

\begin{tcolorbox}[breakable,colback=teal!4,colframe=teal!35!black,title={Before you move on}]
What does \bn{আত্মত্যাগ} mean? (\textquotedblleft{}Self-sacrifice\textquotedblright{}.) Say it once more.
\end{tcolorbox}
\section[\texorpdfstring{\bn{শহিদ} (shôhid)}{শহিদ (shôhid)}]{\bn{শহিদ} (shôhid) — a martyr}

\label{lesson:BN-C317-shohid}



\begin{quote}
Before the new one: say the Bengali for a warrior, then the Bengali for self-sacrifice.
\end{quote}

\subsection*{You'll want to know: \bn{শহিদ}}
\textbf{\bn{শহিদ}} — \emph{shôhid} — \textquotedblleft{}a martyr\textquotedblright{}.

\begin{quote}
\bn{স্বাধীনতার} \bn{যুদ্ধে} \bn{অনেক} \bn{মানুষ} \bn{শহিদ} \bn{হলেন।}
\end{quote}

\begin{quote}
\emph{svādhīnôtār juddhe ônek mānush shôhid hôlen.}
\end{quote}

\begin{quote}
Many people gave their lives in the war of independence.
\end{quote}

\begin{cousinweb}
\textbf{\bn{শহিদ}} (\emph{shôhid}) is usually traced through Persian to Arabic \emph{shahīd}, 'a witness', and so one who bears witness with his life. \bn{শহিদ} \bn{হওয়া}, to die for a cause.
\end{cousinweb}

\begin{grammarlens}[title={What you've built}]
One more word for telling what happened.
\end{grammarlens}

\subsection*{Your turn}
\begin{itemize}
\raggedright
  \item \emph{Say it:} \emph{shôhid}
  \item \emph{Say it:} \emph{shôhid}, once more
  \item \emph{Say it:} \emph{svādhīnôtār juddhe ônek mānush shôhid hôlen.} — Many people gave their lives in the war of independence.
  \item \emph{Recall:} say the Bengali for a warrior, then the Bengali for self-sacrifice, then say \emph{shôhid} again
\end{itemize}

\begin{tcolorbox}[breakable,colback=teal!4,colframe=teal!35!black,title={Before you move on}]
What does \bn{শহিদ} mean? (\textquotedblleft{}A martyr\textquotedblright{}.) Say it once more.
\end{tcolorbox}
\section[\texorpdfstring{\bn{আত্মসমর্পণ} \bn{করা} (ātmôsômôrpôn kôrā)}{আত্মসমর্পণ করা (ātmôsômôrpôn kôrā)}]{\bn{আত্মসমর্পণ} \bn{করা} (ātmôsômôrpôn kôrā) — to surrender}

\label{lesson:BN-C317-atmosomorpon-kora}



\begin{quote}
Before the new one: say the Bengali for self-sacrifice, then the Bengali for a martyr.
\end{quote}

\subsection*{You'll want to know: \bn{আত্মসমর্পণ} \bn{করা}}
\textbf{\bn{আত্মসমর্পণ} \bn{করা}} — \emph{ātmôsômôrpôn kôrā} — \textquotedblleft{}to surrender\textquotedblright{}.

\begin{quote}
\bn{অবশেষে} \bn{শত্রুরা} \bn{আত্মসমর্পণ} \bn{করল।}
\end{quote}

\begin{quote}
\emph{ôbôsheshe shôtrurā ātmôsômôrpôn kôrlô.}
\end{quote}

\begin{quote}
At last the enemy surrendered.
\end{quote}

Past (he/she …): \bn{তিনি} \bn{আত্মসমর্পণ} \bn{করলেন}, \emph{tini ātmôsômôrpôn kôrlen}.

\begin{cousinweb}
\textbf{\bn{আত্মসমর্পণ} \bn{করা}} (\emph{ātmôsômôrpôn kôrā}) puts Sanskrit \emph{ātman}, 'self', and \emph{samarpaṇa}, 'handing completely over' (Monier-Williams), before \bn{করা}: to hand yourself over.
\end{cousinweb}

\begin{grammarlens}[title={What you've built}]
One more word for telling what happened.
\end{grammarlens}

\subsection*{Your turn}
\begin{itemize}
\raggedright
  \item \emph{Say it:} \emph{ātmôsômôrpôn kôrā}
  \item \emph{Say it:} \emph{ātmôsômôrpôn kôrā}, once more
  \item \emph{Say it:} \emph{ôbôsheshe shôtrurā ātmôsômôrpôn kôrlô.} — At last the enemy surrendered.
  \item \emph{Recall:} say the Bengali for self-sacrifice, then the Bengali for a martyr, then say \emph{ātmôsômôrpôn kôrā} again
\end{itemize}

\begin{tcolorbox}[breakable,colback=teal!4,colframe=teal!35!black,title={Before you move on}]
What does \bn{আত্মসমর্পণ} \bn{করা} mean? (\textquotedblleft{}To surrender\textquotedblright{}.) Say it once more.
\end{tcolorbox}
\section[\texorpdfstring{\bn{পরাজয়} (pôrājôy)}{পরাজয় (pôrājôy)}]{\bn{পরাজয়} (pôrājôy) — defeat}

\label{lesson:BN-C317-porajoy}



\begin{quote}
Before the new one: say the Bengali for a martyr, then the Bengali for to surrender.
\end{quote}

\subsection*{You'll want to know: \bn{পরাজয়}}
\textbf{\bn{পরাজয়}} — \emph{pôrājôy} — \textquotedblleft{}defeat\textquotedblright{}.

\begin{quote}
\bn{পরাজয়ের} \bn{পরে} \bn{রাজা} \bn{দেশ} \bn{ছেড়ে} \bn{পালালেন।}
\end{quote}

\begin{quote}
\emph{pôrājôyer pôre rājā desh chheṛe pālālen.}
\end{quote}

\begin{quote}
After the defeat the king fled the country.
\end{quote}

\begin{cousinweb}
\textbf{\bn{পরাজয়}} (\emph{pôrājôy}) is Sanskrit \emph{parājaya}, 'loss, defeat' (Monier-Williams), with \emph{parā-}, 'away', before \emph{jaya}, victory — \textbf{\bn{জয়}}.
\end{cousinweb}

\begin{grammarlens}[title={What you've built}]
One more word for telling what happened.
\end{grammarlens}

\subsection*{Your turn}
\begin{itemize}
\raggedright
  \item \emph{Say it:} \emph{pôrājôy}
  \item \emph{Say it:} \emph{pôrājôy}, once more
  \item \emph{Say it:} \emph{pôrājôyer pôre rājā desh chheṛe pālālen.} — After the defeat the king fled the country.
  \item \emph{Recall:} say the Bengali for a martyr, then the Bengali for to surrender, then say \emph{pôrājôy} again
\end{itemize}

\begin{tcolorbox}[breakable,colback=teal!4,colframe=teal!35!black,title={Before you move on}]
What does \bn{পরাজয়} mean? (\textquotedblleft{}Defeat\textquotedblright{}.) Say it once more.
\end{tcolorbox}
